% Figure: host response timeline from implantation to chronic state.
% Vol VI Ch 12 sec 12.2.
\begin{figure}[htbp]
\centering
\begin{tikzpicture}[
  phase/.style={draw=engblack, minimum height=1.0cm, align=center,
    font=\small},
  label/.style={font=\scriptsize\itshape, anchor=south, align=center},
]

% Time axis
\draw[->, thick] (-0.3, 0) -- (14.5, 0);
\node[font=\scriptsize] at (14.7, -0.3) {time (log)};

% Tick marks and labels
\foreach \x/\t in {0.5/0, 2.5/seconds, 5.0/hours, 7.5/days, 10.0/weeks, 12.5/months--years} {
  \draw (\x, -0.05) -- (\x, 0.05);
  \node[font=\tiny, anchor=north] at (\x, -0.05) {\t};
}

% Phase boxes
\node[phase, fill=engblue!10, minimum width=2.5cm] (adsorb)
  at (2.5, 1.4) {\textbf{Protein adsorption}\\
  (Vroman effect)};
\node[phase, fill=englime!10, minimum width=2.5cm] (acute)
  at (5.5, 1.4) {\textbf{Acute inflammation}\\
  neutrophils};
\node[phase, fill=engred!10, minimum width=2.5cm] (chronic)
  at (8.5, 1.4) {\textbf{Chronic inflammation}\\
  macrophages};
\node[phase, fill=engred!15, minimum width=2.5cm] (granuloma)
  at (11.5, 1.4) {\textbf{Foreign-body reaction}\\
  giant cells, fibrosis};

% Encapsulation (longer arrow extending right)
\node[phase, fill=engblack!10, minimum width=3cm] (capsule)
  at (12.0, 3.0) {\textbf{Fibrous capsule}\\
  (mature; weeks to years)};

% Connecting arrows
\draw[->, thick] (adsorb) -- (acute);
\draw[->, thick] (acute) -- (chronic);
\draw[->, thick] (chronic) -- (granuloma);
\draw[->, thick] (granuloma) -- (capsule);

% Side note: integration paths
\node[label, fill=engblue!5, draw=engblack, minimum width=4cm,
  minimum height=1.5cm, align=center] at (3.0, 3.5)
  {\textbf{Alternative outcomes}\\
  Osseointegration (bone)\\
  Endothelialisation (vasculature)\\
  Degradation (resorbable)};

\end{tikzpicture}
\caption[Host-response timeline]{The canonical foreign-body
response timeline at a biomaterial surface, on a logarithmic time
axis spanning seconds (initial protein adsorption) through months
to years (mature fibrous capsule). Protein adsorption is rapid
and the adsorbed layer is what later cells recognise, not the
nominal material surface; the Vroman effect describes the
sequential displacement of early-adsorbed abundant proteins
(albumin) by lower-abundance higher-affinity proteins (fibrinogen,
factor XII). Acute neutrophilic inflammation peaks within hours;
chronic macrophage-dominated inflammation persists if the implant
is not resorbed; macrophage fusion to foreign-body giant cells
and the surrounding fibrous capsule are the chronic equilibrium
for a non-degradable solid implant. The side box names the
favourable integrative paths available to specific tissue and
device types
\cite{paper:v6c12-anderson-foreign-body-2008,
paper:v6c12-vroman-effect-1969,
text:v6c12-ratner-biomaterials-2020}.}
\label{fig:vol06:ch12:host-response-timeline}
\end{figure}
